\needspace{10\baselineskip}
\section{Arrays e operações numéricas com NumPy}
\label{sec:numpy}

O NumPy é uma biblioteca de computação numérica. Seu objeto central é o \codigo{ndarray}, um array com dimensões, formato e um tipo de dado comum aos elementos \cite{numpy}. Um vetor tem uma dimensão e uma matriz tem duas.

\subsection{Criação de arrays}

\codigo{np.array()} converte uma sequência em array. \codigo{np.zeros()} e \codigo{np.ones()} criam arrays inicializados. Um argumento inteiro define o comprimento de um vetor. Uma tupla pode indicar as dimensões de uma matriz.

\begin{exemplo}
import numpy as np

notas = np.array([8.5, 7.0, 9.2, 5.5])
print(notas.shape)  # (4,)
print(notas.dtype)  # float64

zeros = np.zeros((2, 3), dtype=int)
print(zeros)
# [[0 0 0]
#  [0 0 0]]
print(np.ones(3))  # [1. 1. 1.]
\end{exemplo}

\codigo{np.arange(inicio, fim, passo)} recebe um intervalo e um passo. Para passos inteiros como o deste exemplo, os valores vão até antes de \codigo{fim}. \codigo{np.linspace(inicio, fim, num)} recebe a quantidade de pontos e inclui as extremidades por padrão.

\begin{exemplo}
import numpy as np

print(np.arange(0, 10, 2))     # [0 2 4 6 8]
print(np.linspace(0, 1, 5))   # [0.   0.25 0.5  0.75 1.  ]
\end{exemplo}

\subsection{Formato, índices e eixos}

\codigo{ndim} informa o número de dimensões e \codigo{shape} informa o tamanho em cada uma. Os índices começam em zero. Em uma matriz, a vírgula separa o seletor de linhas do seletor de colunas.

\begin{exemplo}
import numpy as np

matriz = np.array([[1, 2, 3], [4, 5, 6]])
print(matriz.shape)   # (2, 3)
print(matriz.ndim)    # 2
print(matriz[0, 0])   # 1
print(matriz[:, 1])   # [2 5]
print(matriz[1, :])   # [4 5 6]
print(matriz[0, 1:])  # [2 3]
\end{exemplo}

Ao somar ao longo de um eixo, a operação reduz esse eixo. Para uma matriz, \codigo{axis=0} percorre as linhas e retorna uma soma por coluna. \codigo{axis=1} percorre as colunas e retorna uma soma por linha. Sem \codigo{axis}, \codigo{sum()} soma todos os elementos.

\begin{exemplo}
import numpy as np

matriz = np.array([[1, 2, 3], [4, 5, 6]])
print(matriz.sum(axis=0))  # [5 7 9]
print(matriz.sum(axis=1))  # [ 6 15]
print(matriz.sum())        # 21
\end{exemplo}

\codigo{reshape()} reorganiza o formato preservando a quantidade de elementos. Por exemplo, um vetor com seis valores pode ser visto como uma matriz com duas linhas e três colunas.

\begin{exemplo}
import numpy as np

sequencia = np.arange(1, 7)
print(sequencia.reshape(2, 3))
# [[1 2 3]
#  [4 5 6]]
\end{exemplo}

\subsection{Operações vetorizadas e filtros}

Operações vetorizadas aplicam a mesma regra ao array inteiro. Na multiplicação de dois vetores de mesmo tamanho, os elementos de posições correspondentes são multiplicados. Uma comparação gera uma máscara de booleanos, que permite selecionar apenas as posições em que a condição é verdadeira.

\begin{exemplo}
import numpy as np

x = np.array([1, 2, 3])
y = np.array([10, 20, 30])
print(x * y)  # [10 40 90]

notas = np.array([8.5, 7.0, 9.2, 5.5])
print(notas + 0.5)  # [9.  7.5 9.7 6. ]
mascara = notas >= 7
print(mascara)         # [ True  True  True False]
print(notas[mascara])  # [8.5 7.  9.2]
\end{exemplo}

\subsection{Broadcasting}

O \emph{broadcasting} permite operações entre formatos compatíveis. O NumPy compara as dimensões da direita para a esquerda. Em cada posição, os tamanhos precisam ser iguais ou um deles deve ser \(1\). Dimensões à esquerda que estiverem ausentes são tratadas como tendo tamanho \(1\) \cite{numpy}.

\begin{exemplo}
import numpy as np

matriz = np.array([[1, 2, 3], [4, 5, 6]])
pesos = np.array([10, 20, 30])
print(matriz * pesos)
# [[ 10  40  90]
#  [ 40 100 180]]
\end{exemplo}

O vetor tem três valores, a mesma quantidade de colunas da matriz. Por isso, a multiplicação aplica os mesmos três pesos a cada linha. Quando houver erro de compatibilidade, verifique \codigo{shape} para identificar quais dimensões estão sendo comparadas.

\subsection{Transposta e produto de matrizes}
\label{sec:matrizes-numpy}

O NumPy transforma listas retangulares, como as da Seção \ref{sec:matrizes-python}, em arrays de duas dimensões. \codigo{a.T} troca linhas por colunas. Em arrays, \codigo{a * b} multiplica elementos correspondentes, e \codigo{a @ b} calcula o produto matricial \cite{numpy}.

\begin{exemplo}
import numpy as np

a = np.array([[1, 2], [3, 4]])
b = np.array([[5, 6], [7, 8]])

print(a.T.tolist())      # [[1, 3], [2, 4]]
print((a + b).tolist())  # [[6, 8], [10, 12]]
print((a * b).tolist())  # [[5, 12], [21, 32]]
print((a @ b).tolist())  # [[19, 22], [43, 50]]
\end{exemplo}

Para calcular \codigo{a @ b}, a quantidade de colunas de \codigo{a} deve ser igual à quantidade de linhas de \codigo{b}. Cada entrada do resultado combina uma linha de \codigo{a} com uma coluna de \codigo{b}. Por exemplo, o primeiro valor é \(1 \cdot 5 + 2 \cdot 7 = 19\).

\subsection{Funções estatísticas}

\codigo{np.mean()}, \codigo{np.median()}, \codigo{np.min()} e \codigo{np.max()} calculam medidas a partir de um array. Assim como na soma, o parâmetro \codigo{axis} permite obter uma medida por linha ou coluna.

\begin{exemplo}
import numpy as np

notas = np.array([8.5, 7.0, 9.2, 5.5])
print(round(np.mean(notas), 2))   # 7.55
print(np.median(notas))           # 7.75
print(np.min(notas), np.max(notas))  # 5.5 9.2
print(np.percentile(notas, 50))   # 7.75
\end{exemplo}

Se o array contém \codigo{np.nan}, \codigo{np.nanmean()} calcula a média dos valores presentes. Essa escolha mantém o valor ausente fora da conta e precisa ser indicada ao descrever o resultado.

\begin{exemplo}
import numpy as np

notas = np.array([8.5, np.nan, 6.5])
print(np.mean(notas))     # nan
print(np.nanmean(notas))  # 7.5
\end{exemplo}
